% 附录 C：关键结果表
% ---------------------------------------------------------------------------
% 写作阶段填写。把正文中每个问题的关键数值结果集中成表，含不确定性区间，
% 让评阅人不用回翻正文即可核对结论。
% ---------------------------------------------------------------------------
%
% 建议结构（示例，按本题实际情况调整）：
%
% \begin{table}[htbp]
%   \centering
%   \caption{问题二各 BMI 组最佳 NIPT 时点与达标比例（含 Bootstrap 区间）}
%   \begin{tabular}{cccc}
%     \toprule
%     BMI 组 & 推荐时点（周） & Bootstrap 95\% 区间 & 达标比例 \\
%     \midrule
%     $[20,28)$ & 15.2 & $[14.0,16.5]$ & 0.903 \\
%     % ... 按实际结果填写
%     \bottomrule
%   \end{tabular}
% \end{table}
%
% 每个直接答案数值都应在这里给出对应的不确定性区间或灵敏度范围，
% 与正文 \texttt{answerbox} 一一对应。
